\def\level{BTS}  % Classe à droite
\def\course{CIEL 2}
\def\chaptername{Equations différentielles} % Titre au centre
\def\docpart{Solutions des Equations différentielles} % Partie du cours
\def\docname{C211}        % Identifiant à gauche

\pagestyle{cours_FP}
\makeheaderBTS{} % Affiche le bandeau


Dans de nombreuses activités scientifiques, les phénomènes étudiés sont régis par des relations faisant intervenir une fonction et certaines de ses dérivées.

\subsection{En électricité}

\begin{minipage}{8cm}
\includegraphics[scale=0.45]{image2_6.jpg}

\end{minipage}
\begin{minipage}{10cm}
La loi des mailles en électricité nous donne:
$$E=Ri(t)+u(t)$$
L'intensité circulant à l'intérieur d'un condensateur s'exprimant par $i=C\, \dfrac{du}{dt}$, on peut écrire:
$$E=RC\, \dfrac{du(t)}{dt}+u(t)=RC\, u'(t)+u(t)$$
\end{minipage}

\bigskip

\begin{minipage}{8cm}
\includegraphics[scale=0.3]{circuit-tp-cours-2.png}

\end{minipage}
\begin{minipage}{10cm}
La loi des mailles en électricité nous donne:
$$E=u_R+u_L=Ri(t)+u_L(t)$$
La tension aux bornes d'une inductance s'exprimant par $u=L\, \dfrac{di}{dt}$, on peut écrire:
$$E=Ri(t)+L\, \dfrac{di(t)}{dt}+i(t)=Ri(t)+L\, i'\,(t)$$
\end{minipage}

\subsection{En bactériologie}

Soit une  population de bactéries qui se développent dans un milieu. La  \og production \fg{} de bactéries est proportionnelle au nombre de bactéries à chaque instant.\\

On pose $B$ le nombre de bactéries en fonction du temps. On a donc:
$$\dfrac{d B}{d t}=k \times B$$
Avec $B(0)=B_0$ le nombre de bactéries au début de l'expérience et $k$ un coefficient de proportionnalité.

\subsection{En mécanique}


\begin{minipage}{7cm}
\includegraphics[scale=0.35]{2021-07-16_151316.jpg}

\begin{itemize}
	\item P: poids du mobile
	\item R: force de réaction
	\item F: force de rappel
\end{itemize}
\end{minipage}
\begin{minipage}{11cm}
Soit un solide de masse $m$ soumis à un ressort de constante de raideur $k$. On veut connaitre l'évolution de sa position au cours du temps.\\

Selon l'axe $(Ox)$ seule la force de rappel intervient et si on suppose un système sans frottement, on peut écrire le principe fondamental de la dynamique:
$$m \dfrac{d ^2 x}{d t^2} = F_{(Ox)} \Leftrightarrow m \dfrac{d ^2 x}{d t^2} = -kx \Leftrightarrow m \dfrac{d ^2 x}{d t^2} + kx =0$$ 

$$\Leftrightarrow m\, x\, ''(t)+kx(t)=0$$


\end{minipage}

\pagebreak


\begin{definition}[	Equation différentielle]{}
Une \textbf{équation différentielle} est une relation mathématique qui lie une fonction inconnue à une ou plusieurs de ses dérivées. L'inconnue de cette équation n'est pas un nombre mais une fonction. 
\end{definition}

\begin{exemple}[Types d'équations différentielles]{}
Voici quelques exemples d'équations différentielles:
\begin{itemize}
	\item $y' - y = 7$ (équation différentielle du premier ordre à coefficients constants)
	\item $y'' - 3y' + 2y = 0$ (équation différentielle du second ordre à coefficients constants)
	
	\end{itemize}
\end{exemple}


\medskip

\begin{definition}[Solution d'une équation différentielle]{}
On appelle \textbf{solution} d'une équation différentielle $(E)$ sur un intervalle $I$ de $\R$ toute fonction vérifiant la relation $(E)$ sur $I$.

\end{definition}

\medskip

\begin{methode}[Vérification d'une solution]{}
Pour vérifier qu'une fonction $f$ est solution d'une équation différentielle $(E)$, il suffit de:
\begin{enumerate}
	\item Calculer les dérivées de $f$ nécessaires à l'équation $(E)$
	\medskip
	\item Remplacer ces dérivées dans l'équation $(E)$
	\medskip
	\item Simplifier l'expression obtenue et conclure.
\end{enumerate}
\end{methode}

\begin{exemple}[Vérifiez que $g(t)=\e^{-t}$ est solution de $(E):\: y'+y=0$]{}


\begin{blocvisiblex}[height=3.5]
    {\begin{enumerate}
		\item Calcul de la dérivée de $g$: \dotfill{}
		\item Remplacer dans $(E)$: \dotfill{}
		\item Vérification de l'égalité: \dotfill{}
	\end{enumerate}


	}
\begin{enumerate}
	\item Calcul de la dérivée de $g$: $g'(t)=-\e^{-t}$

	\item Remplaçons $g$ et $g'$ dans l'équation $(E)$: $g'(t)+g(t)=-\e^{-t}+\e^{-t}=0$

	\item Vérification de l'égalité:
	L'équation $(E)$ est vérifiée pour tout $t \in \R$. On conclut que $g$ est solution de l'équation différentielle $(E)$.
	
\end{enumerate}

\end{blocvisiblex}


\end{exemple}

\begin{exercice_cours}[Vérification d'une solution]{}
\begin{minipage}{0.30\linewidth}
	Vérifiez que $f(t)=4t^2+t-1$ 
	
	solution de $(E_1): \, y'=8t+1$.

	\medskip

	\dotfill{}

	\medskip{}

	\dotfill{}

	\medskip

	\dotfill{}

	\medskip

	\dotfill{}

	\medskip
	
	\dotfill{}

	\medskip
	
	\dotfill{}

	\medskip
	
	\dotfill{}

	\medskip
	
	\dotfill{}

	\medskip
	
	\dotfill{}

\end{minipage}\hfill{}
\begin{minipage}{0.30\linewidth}
		Vérifiez que $g(x)=4+\e^{-3x}$ 
		
		solution de $(E_2): \, y'+3y=12$.

		
	\medskip

	\dotfill{}

	\medskip{}

	\dotfill{}

	\medskip

	\dotfill{}

	\medskip

	\dotfill{}

	\medskip
	
	\dotfill{}

	\medskip
	
	\dotfill{}

	\medskip
	
	\dotfill{}

	\medskip
	
	\dotfill{}

	\medskip
	
	\dotfill{}
\end{minipage}\hfill{}
\begin{minipage}{0.30\linewidth}
		Vérifiez que $f(x)=\e^{2x}+3$ 
		
		sol. de $(E_3): \, y'-2y=-6$.

		
	\medskip

	\dotfill{}

	\medskip{}

	\dotfill{}

	\medskip

	\dotfill{}

	\medskip

	\dotfill{}

	\medskip
	
	\dotfill{}

	\medskip
	
	\dotfill{}

	\medskip
	
	\dotfill{}

	\medskip
	
	\dotfill{}

	\medskip
	
	\dotfill{}
\end{minipage}
\end{exercice_cours}
